\exerciseeditionsection{Solutions to Wald's Problems}
\ExerciseEditorialNote

\begin{enumerate}[leftmargin=2.8em,itemsep=1.2em]
\WaldSolution{1}
\begin{enumerate}[label=\alph*),leftmargin=2em]
\item Put \(g(r)=\psi(r)f(r,\theta_0)\).  The cutoff makes \(g\) and all its
derivatives vanish at the far end of the ray.  Repeated integration by parts
in
\(\int_0^R r^{k-1}g^{(k)}(r)\dd r\) leaves only the endpoint \(g(0)=f(0)\)
and gives
\[
 f(0)=\frac{(-1)^k}{(k-1)!}
 \int_0^{R(\theta_0)}r^{k-1}g^{(k)}(r)\,\dd r.
\]

\item Average this identity over the cone's solid angle.  Since
\(\dd V=r^{n-1}\dd r\dd\Omega\), the radial weight becomes
\(r^{k-1}\dd r\dd\Omega=r^{k-n}\dd V\).  Thus
\[
 f(0)=C_1\int_Qr^{k-n}\frac{\dd^k}{\dd r^k}(\psi f)\,\dd V,
 \qquad C_1=\frac{(-1)^k}{(k-1)!\Omega}.
\]

\item Cauchy--Schwarz gives a radial factor
\[
 \left(\int_Qr^{2k-2n}\dd V\right)^{1/2},
\]
whose integral near the vertex is
\(\int_0 r^{2k-n-1}\dd r\), finite precisely when \(k>n/2\).  Expanding
the radial derivatives of \(\psi f\) bounds the other factor by a constant
times \(\|f\|_{Q,k}\).  The fixed cone height, angle, and cutoff make the
constant uniform.  Placing such a cone at every point of \(A\) yields
\[
 \sup_{x\in A}|f(x)|\le C\|f\|_{A,k}.
\]
\end{enumerate}

\WaldSolution{2}
\begin{enumerate}[label=\alph*),leftmargin=2em]
\item Decompose
\(F_{ab}=2n_{[a}E_{b]}+\epsilon_{abcd}n^cB^d\).  Contracting
\(\nabla_aF^{ab}=4\pi j^b\) with \(n_b\), and projecting the homogeneous
equation in the same way, gives
\[
 D_aE^a=4\pi\rho,
 \qquad D_aB^a=0.
\]
Extrinsic-curvature terms cancel because \(K_{ab}\) is symmetric while the
relevant spatial tensors are antisymmetric.

\item Choose a global potential \(F_{ab}=2\nabla_{[a}A_{b]}\) and impose
Lorenz gauge.  The equations reduce to the normally hyperbolic system
\[
 \nabla^b\nabla_bA_a-R_a{}^bA_b=0.
\]
The divergence constraints are exactly the compatibility conditions on the
initial data.  Standard global-hyperbolic wave-equation theory supplies a
solution with continuous dependence and the domain-of-dependence property.
The gauge condition propagates because its violation satisfies a homogeneous
scalar wave equation.  Two solutions with the same \(E,B\) data differ by a
pure gauge potential, so their \(F_{ab}\) fields agree.  This proves existence,
uniqueness, and well posedness for \(F_{ab}\).
\end{enumerate}

\WaldSolution{3}
Since \(h_{ab}=g_{ab}+n_an_b\), direct expansion gives
\[
 \frac12\mathcal L_nh_{ab}
 =h_{(a}{}^c\nabla_{|c|}n_{b)}.
\]
Hypersurface orthogonality implies
\(h_a{}^ch_b{}^d\nabla_{[c}n_{d]}=0\), while
\(n^b h_a{}^c\nabla_cn_b=0\).  The projected derivative is therefore
already symmetric and spatial, so
\[
 h_a{}^c\nabla_cn_b=\frac12\mathcal L_nh_{ab}.
\]

\WaldSolution{4}
Using \(K_{ab}=h_a{}^ch_b{}^d\nabla_cn_d\), expand
\[
 D_aK^a{}_b-D_bK
 =h_b{}^c h^{ad}\nabla_a(h_d{}^e\nabla_en_c)
 -h_b{}^c h^{ad}\nabla_c(h_a{}^e\nabla_en_d).
\]
Derivatives of the projectors cancel between the two terms after
\(h_a{}^bn_b=0\) and the symmetry of \(K_{ab}\) are used.  The remainder is
\[
 h_b{}^ch^{ae}(\nabla_a\nabla_e-
 \nabla_e\nabla_a)n_c
 =h_b{}^cR_{cd}n^d,
\]
which is Eq.~\eqref{eq:10.2.24}.

\WaldSolution{5}
At any point choose coordinates adapted to the foliation with
\(g_{00}=-1\), \(g_{0i}=0\), and \(n^\mu=\delta^\mu{}_0\).  In the
second-derivative part of Eq.~\eqref{eq:10.2.26}, the coefficients of every
\(\partial_0^2g_{\alpha\beta}\) cancel in \(G_{00}\) and \(G_{i0}\).
The surviving principal terms contain only spatial second derivatives and
mixed space--time derivatives.  Since the statement is tensorial, it holds in
every adapted coordinate system: \(G_{\mu\nu}n^\nu\) depends on \(h_{ab}\),
\(K_{ab}\), and their spatial derivatives, with no second time derivatives.

\WaldSolution{6}
On a Cauchy surface the gauge-invariant initial data \((\mathbf E,\mathbf B)\)
contain six functions.  The two scalar constraints
\(D_aE^a=D_aB^a=0\) remove two, leaving four freely specifiable phase-space
functions.  A propagating configuration variable and its time derivative
count as a pair, so four phase-space functions correspond to two local
degrees of freedom, the two electromagnetic polarizations.
\end{enumerate}
